\subsection{The remaining ordinary-axis conductor}

\begin{statement}{P4.49}{Proved}
At an ordinary three-axis critical-curve point with uniform
generic transverse type $A_3$ and local defect nine, the
normalization has rational singularities. This includes the
quartic square with exactly one missing coefficient left
by P4.47--P4.48.
\end{statement}
\begin{proof}
Only the one-missing-coefficient case remains. Rescale
its leading conductor form to $x(y+z)$.
As in P4.47, the descent colength one lifts the leading
kernel vector to $g\in I$ with this quadratic part.
Write $g=x(\alpha y+\beta z)+yz\gamma$, where
$\alpha,\beta$ are units and $\gamma(0)=0$.
Replacing $y,z$ by $\alpha y,\beta z$ is an analytic
coordinate change preserving the axes. Put $k=y+z$,
$h=yz$. Then $g=xk-hc$, with $c(0)=0$.
The ideal $I_0=(g,xh,h^2)$ lies in the actual conductor
$I$ and has Hilbert--Burch matrix
\[
                 \begin{pmatrix}h&0\\-k&h\\c&-x\end{pmatrix}.
\]
It has radical $(xy,xz,yz)$, is Cohen--Macaulay, and
agrees with $I$ generically on each axis. The
torsionfree finite-kernel comparison used in P4.47
therefore gives the exact equality $I=I_0$.

There are exactly two required ideal forms. Since
$\mathcal O/(x,k,h)=\mathbb C\{y\}/(y^2)$, analytic
division writes $c=\lambda y+xA+kB+hC$ for a scalar
$\lambda$. Replace $g$ by $g+xhA+h^2C$, and then
set $x'=x-hB$. These generator and coordinate changes
give
\[
                    I=(x'k-\lambda yh,x'h,h^2).
\]
If $\lambda=0$ this is $(xk,xh,h^2)$. If
$\lambda\ne0$, rescaling $x'=\lambda X$ gives
$(Xk-yh,Xh,h^2)$. Every change is invertible near
the origin. We treat the two forms separately.

\paragraph{The ideal $(xk,xh,h^2)$.}
Appending $(a,b,c)$ to its Hilbert--Burch matrix gives
determinant $axk+bxh+ch^2$. The adjugate ring condition
of de Jong--van Straten, Theorem~1.12, applies to the
actual conductor as in P4.41. The colon ideals
\[
 I:h=(x,h),\qquad I:x=(k,h),\qquad I:k=(x,h^2)
\]
give exactly
$f=A x^2k^2+B x^2kh+Cxh^2+Dh^4$, with $A(0)\ne0$.
Complete the square in $k$, absorb $\sqrt A$ into $x$,
and replace $y$ by $y(1+\gamma z)$ to make
$k+\gamma h$ the new $y+z$. These are analytic
coordinate maps; $h$ changes only by a unit.
Absorbing that unit in the coefficients gives
\[
                       f=x^2k^2+xh^2C+h^4D.          \tag{4.49a}
\]
If $C(0)\ne0$, substitute $x=h^2v$ and divide by
$h^4$. The finite $v$-root is simple, giving a smooth
factor $x-h^2v(y,z)$. Its quotient is $xk^2+h^2U$
with $U$ a unit. The coordinate $X=x/U$ turns it
into $Xk^2+h^2$. Its finite birational normalization is
\[
                 u=h/k,\qquad X=-u^2,
                 \qquad u(y+z)=yz.
\]
The last equation is an $A_1$ quadric, so the
normalization has rational singularities.

Suppose $C(0)=0$. Perturb $D$ by a small constant
and the linear part of $C$ generically, retaining
$C(0)=0$, and then replace $C$ by $C-\epsilon$.
On a fixed boundary the transverse tacnode coefficients
are units, so these changes retain boundary $A_3$.
Under $x=\epsilon^3X$, $y=\epsilon Y$,
$z=\epsilon Z$, the leading scaled equation is
\[
 \frac f{\epsilon^8}=F+O(\epsilon),\qquad
 F=X^2K^2+XH^2(L-1)+D_0H^4,
 \quad K=Y+Z,\ H=YZ.
\]
For $L=K$, $D_0=1$, put $W=Y-Z$.
At $W=0$ the $X$-critical coordinate is
$X=-K^2(K-1)/32$, and the remaining derivative is
$K^5(12K^2+7K-3)/512$.
At either root of $12K^2+7K-3$, the full critical
value and Hessian determinant are respectively
\[
 \frac{K^6(K+1)(3K-1)}{1024},\qquad
 -\frac{3K^{10}(K+1)(3K-1)(28K^2+14K-5)}{32768}.
\]
They are nonzero: their numerator resultants with the
critical polynomial are $8748$ and $4923269424$.
Nondegeneracy, nonzero value and avoidance of the axes
hold on a nonempty algebraic open set of coefficients.
An arbitrarily small generic perturbation is therefore
available. The implicit-function theorem lifts these
Morse points through the higher weighted terms, with
nonzero critical values of order $\epsilon^8$.
The boundary comparison embeds their square-$-2$
vanishing spheres in a nearby torus fibre, a
contradiction. Thus $C(0)=0$ is excluded.

\paragraph{The ideal $(g,xh,h^2)$, where $g=xk-yh$.}
Append $(a,b,d)$ to
\[
                         \begin{pmatrix}h&0\\-k&h\\y&-x\end{pmatrix}.
\]
Its determinant is $ag+bxh+dh^2$. The new adjugate
entries give the conditions
\[
       ax,\ ah,\ bx+dh,\ kd+by,\ dh-ay,\ ak+bh\in I.
\]
The exact colon identities are
\[
 I:h=(x,h),\qquad I:x=(xk,h,z^2),\qquad I:x^2=(k,h).
\]
The first two conditions give $a=hA$ modulo $I$;
then $dh-ay$ gives $d=Ay+xD+hE$.
Initially $bx+dh$ permits
$b=-Ak+hB+z^2B_2$ modulo $I$.
The condition $ak+bh$ forces $z^2B_2\in(x,h)$,
so $B_2\in(x,y)$. Its $x$-part lies in $I$, since
$xz^2=h^2-xh$ modulo $I$, and its $y$-part is
absorbed in $hB$. Thus $b=-Ak+hB$ modulo $I$.
The remaining condition is $y(D+B)+kE\in(x,h)$.

Write $T=D+B$. Modulo $(x,h,k)$, the equality
$yT=0$ in $\mathbb C\{y\}/(y^2)$ forces $T(0)=0$.
Write $T=xW+yU+kV$. Since $k$ is a nonzerodivisor
modulo $(x,h)$, the general solution is
$E=-yU-yV+xE_1+hE_2$. Substitution in the determinant
gives $f\in(p)+I^2$, where
\[
                            p=yh^2(x-h).
\]
The exact identities
$xh^3=ygh^2-yp$ and $h^4=ygh^2-kp$ then put the
equation in the form
\[
 f=A g^2+B gxh+Cgh^2+D(xh)^2+\kappa p,
                       \qquad A(0)\ne0.              \tag{4.49b}
\]

We first treat $\kappa(0)\ne0$.
Use the appended coefficients
$a=Ag+Bxh+Ch^2$, $b=Dxh+\kappa yh$,
$d=-\kappa yh$. The actual conductor identifies
$\operatorname{Hom}(I,\mathcal O_S)$ with the
normalization, generated by $1,n_1,n_2$, where the
adjugate fractions are
\[
          n_1=-a/h,\qquad n_2=-(ak+bh)/h^2.
\]
Since $A$ is a unit, the same actual normalization
algebra is generated by $1,v,w$, where
$v=(n_1+Bx+Ch)/A=-g/h$ and $w=n_2$.
Direct multiplication gives the exact relations
\begin{align*}
 Av^2-Bxv-Chv+Dx^2+\kappa xy-\kappa yh&=0,\\
 xw&=Ayv-Bxy-(\kappa+C)yh,\\
 vw&=\kappa h-Dxy,\\
 w^2+[(\kappa+B)y+Ck]w-A\kappa k+ADy^2
             -D(\kappa+C)yh-B\kappa h&=0.
\end{align*}
At the origin they give $v=w=0$, so there is one
normalization point. The last equation eliminates $k$
analytically, with $k$ in the maximal-ideal square.
Thus $h=y(k-y)$ has quadratic initial term $-y^2$.
Write $a=A_0$, $b=B_0$, $d=D_0$, $u=\kappa_0$.
The three quadratic equations become
\[
 av^2-bxv+dx^2+uxy=0,\qquad
 xw-ayv+bxy=0,\qquad vw+uy^2+dxy=0.
\]
Set $V=av-(b/2)x$, $\delta=ad-b^2/4$,
$Y=\delta x+auy$, $W_0=w+(b/2)y$, and
$W=auW_0+\delta V$. This invertible linear change
turns them into
\[
                       V^2+xY=0,\qquad
                       VW+Y^2=0,\qquad xW-YV=0.
\]
They define the irreducible rational normal cubic cone.
The full tangent-cone ideal contains this prime ideal
and has the same dimension two, so it equals it.
Blowing up the normalization point is smooth along its
exceptional rational normal cubic: the two independent
initial equations have independent derivatives along
that smooth projective curve in every blowup chart.
The exceptional curve is $\mathbb P^1$ with normal
bundle $\mathcal O(-3)$. The exact sequences of its
thickenings have successive terms $\mathcal O(3n)$,
all with $H^1=0$. Formal functions gives
$R^1\rho_*\mathcal O=0$. The singularity is rational.

Finally suppose $\kappa(0)=0$. Perturb $D$ and the
linear part of $\kappa$ generically, and replace
$\kappa$ by $\kappa-\epsilon$. The fixed axes retain
boundary $A_3$: equation (4.49b) stays in the same
ring-condition module, and its boundary transverse
quartic coefficients remain units.
Scale $x=\epsilon^2X$, $y=\epsilon Y$, $z=\epsilon Z$.
Put $K=Y+Z$, $H=YZ$, $G=XK-YH$.
After division by $\epsilon^6$, the first terms are
\[
 A_0G^2+\epsilon(B_0GXH+C_0GH^2)
   +\epsilon^2\bigl[D_0(XH)^2+p(L-1)\bigr].
\]
Near a nonzero smooth point of $G=0$, eliminate the
normal Morse coordinate. On that cone the identities
$XH^3=-Yp$, $H^4=-Kp$ absorb the square-completion
corrections into the coefficients, leaving the leading
tangential function
\[
                    \Phi=D'(XH)^2+p(L'-1)|_{G=0}.
\]
Choose $D'=1$, $L'=-Z$ to exhibit a nonempty open
set of suitable coefficients. The parametrization
\[
 K=r,\quad Y=tr,\quad Z=(1-t)r,\quad
 X=t^2(1-t)r^2
\]
gives $\Phi=t^4(1-t)^4r^7[1+r(t^2-t+1)]$.
At $t=1/2$, $r=-7/6$, its value and Hessian are
\[
 -\frac{823543}{573308928},\qquad
 \operatorname{diag}\left(\frac{117649}{1990656},
                          \frac{15647317}{214990848}\right).
\]
The Hessian determinant is
$1840891197733/427972821516288$, nonzero.
The angular critical numerator factors as
$(2t-1)(3t^2-3t-4)$, so this point is simple and
avoids the axes and $K=0$. The generic coefficient
condition is algebraic and nonempty. After normal
elimination, divide the tangential critical equations
by $\epsilon^2$ and apply the implicit-function theorem.
The resulting full Morse points have nonzero critical
values of order $\epsilon^8$ and tend to the origin.
Boundary transport gives the forbidden square-$-2$
vanishing sphere in a nearby torus fibre. Thus
$\kappa(0)=0$ is excluded, completing both ideal forms.
\end{proof}

All adjugates, colon ideals, normalization relations,
linear coordinate identities, and both polar Hessian
calculations are executed in
\texttt{code/part4/a3\_missing\_coefficient.py}.
The saved output reports PASS.

\begin{statement}{P4.50}{Proved}
In the requested full-fibration classes, no support
with positive-genus ruled resolution has positive
Euler number and generic transverse type $A_3$
on every critical-curve component.
\end{statement}
\begin{proof}
P4.42 and P4.44 leave only a contracted horizontal
positive-genus section and a cycle at one three-branch
point of a conductor with two rational normalized
components. The total defect is nine by P4.33 and
(4.42a). P4.48 excludes planar triples; P4R.16
excludes the nonplanar excess-one tacnode-plus-line
case; the exact excess-two classification and conormal
calculation in P4R.17 exclude
every delta-four nonplanar triple. The branchwise bound
$\delta+3+m/2$ excludes all larger delta invariants,
where $m\ge3$ is total multiplicity.
The only remaining triple is the ordinary three-axis
point. P4.48 gives local defect at least nine, so it
uses the entire budget, and P4.47 and P4.49 give
rational normalization there.

Every additional singular critical-curve point has
strictly positive defect and therefore is absent.
At a smooth point of the entire critical curve,
P4R.14 gives the product tacnode or its $A_1$
normalization degeneration; both have rational
normalization. Away from the conductor, P4R.6
excludes exceptional configurations. The normalization
therefore has rational singularities everywhere.
Its minimal resolution cannot contract the horizontal
positive-genus section. This contradiction proves
the assertion.
\end{proof}
